\chapter{Unity}
Unity è un motore di gioco (\textit{game engine})\cite{Unity-1}\cite{Unity-2}\cite{Unity-3} multipiattaforma che consente lo sviluppo di videogiochi e altri contenuti interattivi, come visualizzazioni architettoniche o animazioni 3D in tempo reale. In questo capitolo introduco i concetti di Unity su cui si basa AR Foundation, necessari per comprendere il funzionamento dell'intero sistema.

\section{Pilastri di Unity}
Unity si basa su tre concetti. Il \textbf{GameObject} è l'entità di base della scena (un oggetto 3D, una luce, la camera): di per sé è un contenitore vuoto, il cui comportamento è definito dai \textit{componenti} associati. Il \textbf{Transform} è il componente presente in ogni GameObject e gestisce posizione (x, y, z), rotazione e scala nello spazio 3D della scena. I \textbf{MonoBehaviour} sono gli script C\# che si agganciano a un GameObject come componenti; ne fanno parte anche quelli di AR Foundation, come \texttt{ARPlaneManager} e \texttt{ARTrackedImageManager}. Questi script sfruttano i metodi del ciclo di vita di Unity (\texttt{Awake()}, \texttt{OnEnable()}, \texttt{Update()}, \texttt{OnDisable()}), richiamati automaticamente dal motore, ad esempio per aggiornare a ogni fotogramma gli oggetti 3D in base al movimento dello smartphone.

\section{Camera}
In un progetto senza AR la Camera è l'occhio del giocatore nello spazio virtuale. In AR svolge invece un doppio compito:

\begin{itemize}
    \item \textbf{Movimento}: la camera virtuale deve muoversi nella scena 3D esattamente come il telefono si muove nella stanza reale. Il componente \texttt{ARPoseDriver} aggiorna il Transform della Main Camera con la posa (posizione e rotazione) stimata dallo SLAM (\textit{Simultaneous Localization and Mapping}), ovvero la stima simultanea della posizione del dispositivo e della mappa dell'ambiente circostante.
    \item \textbf{Sfondo}: normalmente lo sfondo è un colore pieno o un cielo (\textit{Skybox}). Tramite \texttt{ARCameraBackground}, AR Foundation lo sostituisce con il flusso video in tempo reale della fotocamera del telefono.
\end{itemize}

\section{Sistema di coordinate}
Unity usa uno spazio di coordinate cartesiane 3D con origine in $(0,0,0)$, in cui per convenzione un'unità corrisponde a un metro. ARCore e ARKit creano invece un proprio sistema di coordinate (\textit{session space}), la cui origine coincide con la posizione del dispositivo all'avvio della sessione. In Unity, quando un GameObject è posto dentro un altro (relazione padre-figlio), le coordinate del figlio sono relative al padre. AR Foundation sfrutta questo meccanismo con un oggetto padre della camera (\texttt{XR Origin}), che converte il session space nel sistema di coordinate del mondo di Unity.

\section{Prefab e istanziazione}
Un \textit{Prefab} è un oggetto salvato come risorsa riutilizzabile, da cui Unity crea copie (istanze) a runtime con la funzione \texttt{Instantiate()}. Quando il telefono rileva un piano fisico (un tavolo, un pavimento), AR Foundation non disegna la geometria dal nulla: \texttt{ARPlaneManager} istanzia il prefab assegnato per ogni piano rilevato, aggiornandolo o rimuovendolo al variare del rilevamento. Lo stesso vale per ancore e altri \textit{trackable}. Lo sviluppatore usa invece \texttt{Instantiate()} per posizionare nella scena i propri contenuti virtuali.

\section{Raycasting}
Per capire ``dove tocca il dito sullo schermo nel mondo reale'' si usa il \textit{raycasting}, cioè il lancio di un raggio invisibile dalla camera attraverso un punto dello schermo, per trovare ciò che incontra lungo il suo percorso:

\begin{itemize}
    \item \textbf{Physics.Raycast}: è il raycasting standard di Unity. Il raggio (ottenibile con \texttt{Camera.ScreenPointToRay}) è testato contro i collider degli oggetti virtuali e restituisce il primo oggetto colpito.
    \item \textbf{ARRaycastManager}: applica lo stesso concetto, ma non usa il motore fisico di Unity: il raggio è testato contro i piani e i punti caratteristici (\textit{feature point}) rilevati dal sistema AR. Il metodo \texttt{Raycast} restituisce una lista di \texttt{ARRaycastHit} ordinata per distanza, ciascuno con una \texttt{pose} (posizione e rotazione del punto colpito); il tipo di superfici considerate si sceglie con \texttt{TrackableType}.
\end{itemize}

La posa restituita da \texttt{ARRaycastManager} viene tipicamente usata per istanziare un prefab sulla superficie toccata dall'utente.